\begin{proof}
Put $g(u)=m(u)-qu$ and $F(u)=\min_{u\le v\le1}g(v)$.
The function $F$ is nondecreasing and Lipschitz: if $u<v$, a minimizer
for $F(u)$ either lies in $[v,1]$, or comparison with $g(v)$ bounds
$F(v)-F(u)$ by $\Lip(g)(v-u)$. Moreover, $F(0)=0$ and
$F(1)\ge s-q>t-q$. Absolute continuity gives an interior
differentiability point $A$ at which $F'(A)>t-q$.
At this point $F(A)=g(A)$, since a strict gap would force the minimum
to be attained a fixed distance to the right and would make $F$ locally
constant. There are contact points $B>A$ arbitrarily close to $A$:
otherwise $F$ would be constant immediately to the right of $A$.
Choose such a $B<\min(2A,1)$ so close that
$F(w)-F(A)\ge(t-q)(w-A)$ on $[A,B]$.
Using $g(w)\ge F(w)$ proves Equation~\eqref{prep:contact-middle}.
The first inequality follows from $m(w)\ge sw$, and the last follows
from $g(w)\ge F(B)=g(B)$ for $w\ge B$.
\end{proof}
